First Dharma concerns the relation samapatti-pratipatti
and the idea of Substantial Knowledge.
There are three species of Substance:
Immutable, Intelligible, Perceptible
What is intelligible substance?

Prajna: Samapatti concernes itself with divisions and principles (prajnas)
Dharma: which is not a Being of Science
in that it is Science as such...a principle as such and a principle-providing.
so it places divisions in beings and is responsible for their genesis
and nongenesis. Only beings which have genesis and nongenesis are definable.

I.41
ksina-vrtter abhijatasyeva maner grahitr-grahana-grahyesu

tat-stha-tad-anjanata samapatti

I.42
tatra sabdartha-jnana-vikalpa sankirna savitarka samapatti

I.43
smrti-parisuddhau svarupa-sunyevartha-matra-nirbhasa nirvitarka

I.44
etayaiva savicara nirvicara ca suksma-visaya vyakhyata

I.45
suksma-visayatvam calinga-paryavasanam

I.46
ta eva sabija samadhi

one first needs to understand
bija-sabija as genus-species
bheda as difference
eka-rupatva , genesis of form
essentially we are forming sciences here
and I think Yoga is written by Guru of Gurus for Gurus.
I dont think it is for ordinary people at all any more.

1. [It is a disjunction into] principles,
with each being equally a principle
of unity and of disjunction.

2. [It provides a deduction from this to a general]
schema of the total empirical domain
according to the form of its genetic principle.

All logical ideas have three moments.

The relation samapatti-pratipatti has a foundational being

a disjunct into principles:
vitarka-vicara

with each being equally a principle of unity and disjunction
vitarka: savitarka-nirvitarka
savitarka: principle of disjunction: ordinary conceiving
nirvitarka: principle of unity: genetic conceiving
vicara: savicara-nirvicara
savicara: principle of disjunction. ordinary conceiving
nirvicara: principle of unity: genetic conceiving

now this should be seen as four locations of the inner instrument
from which pratyaya are accessed,

a samskara is a species of pratyaya,
a pratyaya of Substantial knowledge
these are four objective principles

genetic conceiving is artha-matra-nirbhasa
artha-matra-nirbhasa is speculative reasoning

ordinary conceiving is dhyana.
it contains errors because it is involved with sabda-artha-jnana-vipalka
once clarity is achieved. this extension to the science has to be
integrated into the whole.

all logical ideas contain three moments but in a divided concept,
which all sciences are, the idea of the parts is replaced by an
"aspect" of the idea of the whole. so this self-enclosure
replaces the moment of the idea of the part with an aspect
of the idea of the whole.
so the two parts lose a moment essentially and is transferred
to the speculative moment of the whole.
